
% !TEX program = xelatex
% !TEX encoding = UTF-8
\documentclass[12pt,fontset=none]{ctexart}

\usepackage[margin=2cm,includefoot]{geometry}
\usepackage{fontspec}
\usepackage{amsmath,amssymb}
\usepackage{enumitem}
\usepackage{booktabs}
\usepackage{array}
\usepackage[hidelinks]{hyperref}

% Fonts
\setmainfont{Palatino}[Ligatures=TeX]
\setCJKmainfont{Songti SC}[
  Script = Default,
  BoldFont = Songti SC Bold,
  ItalicFont = Songti SC,
  SlantedFont = Songti SC,
  BoldItalicFont = Songti SC Bold,
  BoldSlantedFont = Songti SC,
]
\setCJKsansfont{Songti SC}[
  Script = Default,
  BoldFont = Songti SC Bold,
]
\setCJKmonofont{Songti SC}[Script=Default]

\ctexset{
  section = {
    format = \Large\bfseries\raggedright,
    aftername = \quad,
  },
  subsection = {
    format = \large\bfseries\raggedright,
    aftername = \quad,
  },
  subsubsection = {
    format = \normalsize\bfseries\raggedright,
    aftername = \quad,
  },
}

\pagestyle{plain}

\setlist[itemize]{
  leftmargin=1.5em,
  itemsep=0.15em,
  topsep=0.25em
}
\setlist[enumerate]{
  leftmargin=1.6em,
  itemsep=0.2em,
  topsep=0.3em
}
\renewcommand{\arraystretch}{1.25}

\newcommand{\coursetitle}[3]{%
  \begin{center}
    {\LARGE #1\par}
    \vspace{1.4em}
    {\large #2\par}
    \vspace{0.3em}
    {\large #3\par}
  \end{center}
}

\hypersetup{
  pdftitle={Signals in the Noise:
  An Introduction to Signal Processing}
}

\begin{document}

% ========================================
% English
% ========================================

\renewcommand{\theHsection}{en.\arabic{section}}

\coursetitle
  {\textbf{Signals in the Noise: \\
  An Introduction to Signal Processing}}
  {Tutor: Qianshuo Ye}
  {Winter 2026}

\thispagestyle{plain}

\section{Course Information}

\begin{tabular}{@{}>{\bfseries}l l@{}}
Class hours & 12 (6 $\times$ 2-hour sessions) \\
Meeting time & TBD \\
Format & Lectures and hands-on practicals \\
Delivery mode & Online (live, interactive sessions) \\
Tools & Google Colab, Python, Audacity (optional)
\end{tabular}

\section{Course Description}

How can we recover meaningful information from noisy
recordings? This course introduces the mathematical
foundations of signal processing through audio analysis
and noise reduction. Students will explore how sounds
are represented digitally, how Fourier transforms reveal
their frequency components, and how convolution and
digital filtering can be used to process signals.

The course also introduces fundamental ideas in
statistical signal processing, including random noise,
probability distributions, correlation, and signal
estimation. Through hands-on exercises using Audacity
and Python, students will apply these concepts to a
mini-project in which they recover hidden messages
from noisy audio recordings.

\section{Prerequisites}

The course is designed for high school students
with a basic understanding of algebra, functions,
and trigonometry. Familiarity with elementary
probability and statistics is helpful but not required.
No prior knowledge of signal processing is assumed.

Basic Python programming experience is desirable. Students can familiarise themselves with Python by completing the course \href{https://github.com/CambridgeEngineering/PartIA-Computing-Michaelmas}{PartIA-Computing-Michaelmas}.

\section{Learning Objectives}

By the end of this course, students will be able to:

\begin{enumerate}
  \item Explain how audio signals are represented and
  sampled, and distinguish between time-domain and
  frequency-domain representations.

  \item Understand the principles of Fourier transforms,
  convolution, and digital filtering, and apply them
  to simple audio signals.

  \item Describe noise using basic statistical concepts,
  including probability distributions, variance,
  correlation, and signal-to-noise ratio (SNR).

  \item Implement basic audio noise reduction methods
  using Python and evaluate their performance.

  \item Apply signal processing techniques to recover
  hidden information from noisy audio recordings.
\end{enumerate}

\section{Syllabus}

\subsection{Signals, Sound and Sampling}

\begin{itemize}
  \item Sound waves: amplitude, frequency, period,
  and phase.
  \item Continuous-time and discrete-time signals.
  \item Superposition of sinusoidal signals.
  \item Digital sampling, sampling rate, Nyquist
  frequency, and aliasing.
\end{itemize}

Key equations:
\begin{align*}
x(t) &= A\sin(2\pi ft+\phi),\\
x[n] &= x(n/f_s),\\
f_s &> 2f_{\max}.
\end{align*}

\subsection{Fourier Analysis}

\begin{itemize}
  \item Time-domain and frequency-domain representations.
  \item Fourier transform and frequency decomposition.
  \item Discrete Fourier transform (DFT) and fast
  Fourier transform (FFT).
  \item Magnitude spectrum, phase spectrum,
  and frequency resolution.
\end{itemize}

Key equations:
\begin{align*}
X[k] &= \sum_{n=0}^{N-1}
x[n]e^{-i2\pi kn/N},\\
f_k &= \frac{kf_s}{N},\qquad
\Delta f = \frac{f_s}{N}.
\end{align*}

\subsection{Convolution and Digital Filtering}

\begin{itemize}
  \item Linear time-invariant (LTI) systems
  and impulse responses.
  \item Discrete convolution and moving average filters.
  \item Low-pass, high-pass, and band-pass filtering.
  \item Frequency responses and the convolution theorem.
\end{itemize}

Key equations:
\begin{align*}
y[n] &= \sum_k x[k]h[n-k],\\
y[n] &= \frac{1}{M}
\sum_{k=0}^{M-1}x[n-k],\\
Y(\omega) &= X(\omega)H(\omega).
\end{align*}

\subsection{Statistical Signal Processing}

\begin{itemize}
  \item Deterministic and random signals.
  \item Random variables and Gaussian noise.
  \item Mean, variance, and standard deviation.
  \item White noise, coloured noise, and correlation.
  \item Signal power and signal-to-noise ratio (SNR).
\end{itemize}

Key equations:
\begin{align*}
y[n] &= x[n]+\epsilon[n],\\
\epsilon[n] &\sim \mathcal{N}(0,\sigma^2),\\
\sigma^2 &= \mathbb{E}[(X-\mu)^2],\\
R_x[\tau] &= \mathbb{E}[x[n]x[n-\tau]],\\
\mathrm{SNR}_{\mathrm{dB}}
&=10\log_{10}\frac{P_x}{P_\epsilon}.
\end{align*}

\subsection{Audio Noise Reduction}

\begin{itemize}
  \item Noise reduction as a signal estimation problem.
  \item Time-domain and frequency-domain filtering.
  \item Introduction to spectrograms and
  short-time Fourier transforms (STFT).
  \item Noise power estimation and basic Wiener filtering.
  \item Evaluation using mean squared error (MSE)
  and SNR improvement.
\end{itemize}

Key equations:
\begin{align*}
\hat X[k] &= H[k]Y[k],\\
H[k] &= \frac{P_x[k]}{P_x[k]+P_\epsilon[k]},\\
\mathrm{MSE}
&= \frac{1}{N}\sum_{n=0}^{N-1}
(x[n]-\hat x[n])^2.
\end{align*}

\subsection{Signal Recovery Mini-Project}

\begin{itemize}
  \item Identifying noise characteristics from
  audio waveforms and frequency spectra.
  \item Selecting and implementing suitable
  noise reduction methods.
  \item Comparing performance under different
  noise types and SNR levels.
  \item Recovering and decoding hidden messages.
  \item Interpreting results and discussing
  the limitations of different methods.
\end{itemize}

\section{Schedule}

\begin{center}
\small
\begin{tabular}{@{}c p{4cm} p{10.8cm}@{}}
\toprule
Session & Topic & Practical Exercise \\
\midrule

1 & Signals, Sound and Sampling
  & Record and visualise audio in Audacity;
    generate sinusoidal signals in Python
    and investigate the effects of sampling. \\

2 & Fourier Analysis
  & Analyse audio spectra using Audacity
    and Python FFT; identify hidden frequency
    components in a mixed signal. \\

3 & Convolution and Digital Filtering
  & Implement a moving average filter;
    compare low-pass and high-pass filtering
    on audio recordings. \\

4 & Statistical Signal Processing
  & Generate different types of noise;
    examine their distributions and spectra;
    add noise at specified SNR levels. \\

5 & Audio Noise Reduction
  & Implement and compare basic denoising
    methods; evaluate their performance
    using MSE and SNR improvement. \\

6 & Signal Recovery Mini-Project
  & Recover hidden messages from noisy
    audio recordings and compare the
    effectiveness of different methods. \\

\bottomrule
\end{tabular}
\end{center}

\section{Deliverables}

Students are expected to complete a mini-project
and submit:

\begin{enumerate}
  \item A Python Jupyter Notebook implementing
  audio noise reduction methods and evaluating
  their performance on noisy recordings.

  \item A brief project report (2--3 pages)
  describing the methods, results, and key
  observations from the signal recovery challenge.
\end{enumerate}

\section{Assessment}

The course will be assessed through a practical
signal recovery challenge. Students will be given
audio recordings contaminated by different types
and levels of noise and tasked with recovering
the hidden messages.

Assessment will consider the accuracy of message
recovery, the effectiveness of the signal processing
methods, and the clarity of the submitted explanations.

% ========================================
% Chinese
% ========================================

\clearpage
\setcounter{section}{0}
\renewcommand{\theHsection}{zh.\arabic{section}}

\coursetitle
  {\textbf{噪声中的信号：信号处理导论}}
  {教师：叶千硕}
  {冬季2026}

\thispagestyle{plain}

\section{课程信息}

\begin{tabular}{@{}>{\bfseries}l l@{}}
总课时 & 12小时（6次课，每次2小时） \\
上课时间 & 待定 \\
授课形式 & 理论讲解与实践操作相结合 \\
授课平台 & 线上（直播，互动式教学） \\
授课工具 & Google Colab, Python, Audacity（可选） \\
\end{tabular}

\section{课程简介}

如何从充满噪声的录音中恢复有用的信息？
本课程以音频信号为切入点，介绍信号处理的
基本数学原理及其实际应用。
学生将学习声音的数字化表示、
傅里叶变换与频谱分析、卷积与数字滤波等核心概念。

课程还将介绍统计信号处理的基本思想，
包括随机噪声、概率分布、相关性与信号估计。
通过 Audacity 和 Python 实践，
学生将运用所学知识完成一项小型项目：
从受到噪声干扰的音频中恢复隐藏的信息。

\section{先修要求}

本课程面向具备基本高中数学知识的学生，
建议掌握代数、函数和三角函数的基础概念。
具备初步的概率与统计知识将有所帮助，
但不要求事先学习过信号处理。

建议学生具备基本的 Python 编程经验。
学生可以通过完成课程 \href{https://github.com/CambridgeEngineering/PartIA-Computing-Michaelmas}{PartIA-Computing-Michaelmas} 中的练习来熟悉 Python。
实践环节将提供必要的示例代码与指导。

\section{学习目标}

学完本课程后，学生应能够：

\begin{enumerate}
  \item 理解音频信号的基本表示与采样方法，
  并区分时域与频域表示。

  \item 理解傅里叶变换、卷积和数字滤波的基本原理，
  并将其应用于简单音频信号。

  \item 使用概率分布、方差、相关性和信噪比等
  基本统计概念描述噪声。

  \item 使用 Python 实现基础音频降噪方法，
  并评价其处理效果。

  \item 运用信号处理方法从含噪音频中
  恢复隐藏的信息。
\end{enumerate}

\section{教学大纲}

\subsection{信号、声音与采样}

\begin{itemize}
  \item 声波的基本性质：振幅、频率、周期与相位。
  \item 连续时间信号与离散时间信号。
  \item 正弦信号的叠加原理。
  \item 数字采样、采样率、奈奎斯特频率与混叠现象。
\end{itemize}

核心公式：
\begin{align*}
x(t) &= A\sin(2\pi ft+\phi),\\
x[n] &= x(n/f_s),\\
f_s &> 2f_{\max}.
\end{align*}

\subsection{傅里叶分析}

\begin{itemize}
  \item 信号的时域与频域表示。
  \item 傅里叶变换与频率分解的基本思想。
  \item 离散傅里叶变换（DFT）
  与快速傅里叶变换（FFT）。
  \item 幅度谱、相位谱与频率分辨率。
\end{itemize}

核心公式：
\begin{align*}
X[k] &= \sum_{n=0}^{N-1}
x[n]e^{-i2\pi kn/N},\\
f_k &= \frac{kf_s}{N},\qquad
\Delta f = \frac{f_s}{N}.
\end{align*}

\subsection{卷积与数字滤波}

\begin{itemize}
  \item 线性时不变系统（LTI）与冲激响应。
  \item 离散卷积与移动平均滤波。
  \item 低通、高通与带通滤波器。
  \item 频率响应与卷积定理。
\end{itemize}

核心公式：
\begin{align*}
y[n] &= \sum_k x[k]h[n-k],\\
y[n] &= \frac{1}{M}
\sum_{k=0}^{M-1}x[n-k],\\
Y(\omega) &= X(\omega)H(\omega).
\end{align*}

\subsection{统计信号处理}

\begin{itemize}
  \item 确定性信号与随机信号。
  \item 随机变量与高斯噪声。
  \item 均值、方差与标准差。
  \item 白噪声、有色噪声及相关性。
  \item 信号功率与信噪比（SNR）。
\end{itemize}

核心公式：
\begin{align*}
y[n] &= x[n]+\epsilon[n],\\
\epsilon[n] &\sim \mathcal{N}(0,\sigma^2),\\
\sigma^2 &= \mathbb{E}[(X-\mu)^2],\\
R_x[\tau] &= \mathbb{E}[x[n]x[n-\tau]],\\
\mathrm{SNR}_{\mathrm{dB}}
&=10\log_{10}\frac{P_x}{P_\epsilon}.
\end{align*}

\subsection{音频降噪}

\begin{itemize}
  \item 将降噪理解为信号估计问题。
  \item 时域滤波与频域滤波。
  \item 声谱图与短时傅里叶变换（STFT）简介。
  \item 噪声功率估计与维纳滤波基础。
  \item 使用均方误差（MSE）
  和信噪比改善量评价降噪效果。
\end{itemize}

核心公式：
\begin{align*}
\hat X[k] &= H[k]Y[k],\\
H[k] &= \frac{P_x[k]}{P_x[k]+P_\epsilon[k]},\\
\mathrm{MSE}
&= \frac{1}{N}\sum_{n=0}^{N-1}
(x[n]-\hat x[n])^2.
\end{align*}

\subsection{信号恢复小项目}

\begin{itemize}
  \item 通过音频波形与频谱识别噪声特征。
  \item 选择并实现合适的降噪方法。
  \item 比较不同噪声类型和信噪比下的处理效果。
  \item 从含噪录音中恢复并破译隐藏的信息。
  \item 分析实验结果及不同方法的局限性。
\end{itemize}

\section{教学进度}

\begin{center}
\small
\begin{tabular}{@{}c p{3.5cm} p{11.3cm}@{}}
\toprule
课次 & 主题 & 实践内容 \\
\midrule

1 & 信号、声音与采样
  & 使用 Audacity 录制并观察音频波形；
    使用 Python 生成正弦信号，
    探索不同采样率的影响。 \\

2 & 傅里叶分析
  & 使用 Audacity 和 Python FFT 分析音频频谱，
    识别混合信号中的隐藏频率。 \\

3 & 卷积与数字滤波
  & 实现移动平均滤波器；
    比较低通和高通滤波对音频信号的影响。 \\

4 & 统计信号处理
  & 生成不同类型的噪声，
    比较其概率分布与频谱，
    并向录音中加入指定信噪比的噪声。 \\

5 & 音频降噪
  & 实现并比较基础降噪方法；
    使用均方误差与信噪比改善量评价效果。 \\

6 & 信号恢复小项目
  & 从受到噪声污染的录音中恢复隐藏信息，
    并比较不同处理方法的有效性。 \\

\bottomrule
\end{tabular}
\end{center}

\section{课程成果}

课程结束时，学生需要完成一项小型项目，
并提交以下成果：

\begin{enumerate}
  \item 一份 Python Jupyter Notebook，
  实现基础音频降噪方法，
  并在含噪录音上评价其效果。

  \item 一份简短的项目报告（2--3页），
  介绍信号恢复挑战中使用的方法、
  处理结果及主要观察。
\end{enumerate}

\section{课程考核}

课程以实践性的信号恢复挑战作为主要考核形式。
学生将获得受到不同类型、
不同强度噪声干扰的音频，
并运用课堂所学知识，
选择合适的方法恢复其中隐藏的信息。

考核将综合考虑信息恢复的准确率、
降噪方法的有效性，
以及对处理过程与结果解释的清晰程度。

\end{document}